% ============================================================
%  part5/ch29c-mso-interlude.tex
%  Interlude: Mereological Space Ontology
%  (short chapter between entropy-causality and observer-horizons)
% ============================================================

\chapter{Interlude: Mereological Space Ontology}
\label{ch:mso}

\begin{WhyChapter}
  Chapter~\ref{ch:analysis-synthesis} established the adjunction
  $F \dashv G$ between decomposition and reconstruction.
  Chapter~\ref{ch:monoturn} will identify the Monoturn as the
  fixed point.
  This short interlude makes explicit the \emph{physical
  interpretation} of that adjunction: the forward and reverse
  Mereological Space Ontology (MSO).
  Without this bridge, readers may absorb the category theory
  without recognising that $F$ and $G$ are describing how
  reality decomposes into parts and how observers reconstruct
  wholes from parts.
\end{WhyChapter}


\begin{ChapterIntro}
There are two questions one can ask about any system.

What is it made of? And what larger structure does it belong to?

The first question leads toward decomposition. The second leads toward reconstruction. Most scientific theories focus heavily on the first. They seek smaller constituents, finer scales, deeper layers. Yet reconstruction is equally important. Observers never encounter isolated fragments. They encounter fragments embedded within larger wholes.

The Mereological Space Ontology treats decomposition and reconstruction as complementary operations acting on the same underlying reality. The resulting picture is neither reductionist nor holistic. It is relational.
\end{ChapterIntro}

\section{The forward MSO: decomposition}

The forward MSO is the chain of decompositions through which
the Monoturn becomes accessible to local observation:

\begin{equation}
  \begin{aligned}
  \mathcal{M}
  &\xrightarrow{F_1} \text{Superclusters}
   \xrightarrow{F_2} \text{Filaments}
   \xrightarrow{F_3} \text{Galaxies}\\
  &\xrightarrow{F_4} \text{Planets}
   \xrightarrow{F_5} \text{Life}
   \xrightarrow{F_6} \text{Cells}\\
  &\xrightarrow{F_7} \text{Atoms}
   \xrightarrow{F_8} \text{Fields}
   \xrightarrow{F_9} \text{Observer.}
  \end{aligned}
  \label{eq:forward-mso}
\end{equation}

Each arrow $F_i$ is a decomposition functor of the type
developed in Chapter~\ref{ch:analysis-synthesis}.
Each arrow loses information: $\ker(F_i) \neq 0$.

\begin{proposition}[Forward MSO loses information at every step]
  \label{prop:fmso-loss}
  \statusproven{}
  For each arrow $F_i$ in the forward MSO,
  $\ker(F_i) \neq 0$.
  The composition $F = F_9 \circ \cdots \circ F_1$ satisfies
  $\ker(F) \supsetneq \ker(F_1)$.
\end{proposition}

\begin{proof}
  Each $F_i$ maps a richer structure to a coarser one.
  By Theorem~\ref{thm:proj-composition}, kernels accumulate:
  $\ker(F_j \circ F_i) \supseteq \ker(F_i)$.
  The containment is strict because each level introduces
  new equivalences.
\end{proof}

\section{The reverse MSO: reconstruction}

The reverse MSO traces the path from observer back toward
the Monoturn:

\begin{equation}
  \begin{aligned}
  \text{Observer}
  &\xrightarrow{G_9} \text{Fields}
   \xrightarrow{G_8} \text{Atoms}
   \xrightarrow{G_7} \text{Cells}\\
  &\xrightarrow{G_6} \text{Life}
   \xrightarrow{G_5} \text{Planets}
   \xrightarrow{G_4} \text{Galaxies}\\
  &\xrightarrow{G_3} \text{Filaments}
   \xrightarrow{G_2} \text{Superclusters}
   \xrightarrow{G_1} \widehat{\mathcal{M}}.
  \end{aligned}
  \label{eq:reverse-mso}
\end{equation}

Each $G_i$ is a reconstruction functor.
The composition $G = G_1 \circ \cdots \circ G_9$ gives
$\widehat{\mathcal{M}} = G(F(\mathcal{M}))$.

\begin{proposition}[Forward and reverse MSO are not inverses]
  \label{prop:mso-not-inverse}
  \statusproven{}
  $G \circ F \neq \mathrm{id}_{\mathbf{Glob}}$ in general.
  Specifically, $\widehat{\mathcal{M}} \subsetneq \mathcal{M}$.
\end{proposition}

\begin{proof}
  By Theorem~\ref{thm:analysis-synthesis-differ}, $G \circ F$
  is an isomorphism only if $F$ is fully faithful, which fails
  when $\ker(F) \neq 0$.
  Since information is lost at every step of the forward MSO,
  the reconstruction $\widehat{\mathcal{M}}$ cannot recover
  the full $\mathcal{M}$.
\end{proof}

\section{Forward and reverse MSO as adjoint operations}

\begin{theorem}[MSO adjunction]
  \label{thm:mso-adjunction}
  \statusproven{}
  The forward MSO $F$ and reverse MSO $G$ form an adjunction:
  \[
    F \dashv G,
  \]
  with unit $\eta : \mathrm{id} \to G \circ F$ and counit
  $\varepsilon : F \circ G \to \mathrm{id}$.
\end{theorem}

\begin{proof}
  This is Theorem~\ref{thm:adjunction} applied to the specific
  case where $F$ is the forward MSO decomposition chain and
  $G$ is the reverse MSO reconstruction chain.
\end{proof}

\begin{InterpBox}[What the adjunction means for part-whole relations]
  The forward MSO asks: \emph{what is this made of?}
  The reverse MSO asks: \emph{what must this be part of?}
  These are not the same question.
  The adjunction $F \dashv G$ is the mathematical statement
  that they are \emph{dual} questions: maps from decompositions
  to local observations correspond to maps from global
  reconstructions to wholes.
  The Monoturn is the object that answers both questions
  simultaneously --- and the answer is always incomplete,
  because $\ker(F) \neq 0$.
\end{InterpBox}

\section{The MSO as a resolution of the Monoturn}

\begin{definition}[MSO resolution]
  The forward MSO chain~\eqref{eq:forward-mso} is a
  \emph{resolution} of the Monoturn: a sequence of
  progressively coarser approximations, each accessible
  to observers at the corresponding scale.
\end{definition}

\begin{proposition}[Observers are located within the MSO]
  \label{prop:observers-in-mso}
  An observer at scale $k$ (e.g., a biologist at the cell scale)
  has direct access to $F_k(\mathcal{M})$ and must reconstruct
  both finer scales ($G_{k+1} \circ \cdots$) and coarser scales
  ($F_{k-1} \circ \cdots$) indirectly.
\end{proposition}

\begin{WorkedExample}[A cosmologist's MSO position]
  A cosmologist works at the galaxy/supercluster scale:
  $F_3(\mathcal{M})$.
  They reconstruct finer scales (individual stars, planets,
  life) indirectly through theory.
  They reconstruct coarser scales (the full Monoturn,
  beyond-horizon structure) through the admissibility sheaf
  and face $H^1(\mathcal{M},\mathcal{A}) \neq 0$ as an
  obstruction.
  Their horizon-relative model
  $\widehat{\mathcal{M}}_O = G(F(\mathcal{M}))$
  is a partial reconstruction at their MSO scale.
\end{WorkedExample}

\begin{ForwardRef}
  The Monoturn (Chapter~\ref{ch:monoturn}) is the global object
  that the reverse MSO attempts but cannot fully reconstruct.
  Observer horizons (Chapter~\ref{ch:observer-horizons})
  are the mechanism that limits how far up the reverse MSO
  chain any single observer can reach.
  Sheaf cohomology (Chapter~\ref{ch:sheaf}) measures the
  failure of that reconstruction.
\end{ForwardRef}

\section{Horizon non-uniqueness}

\begin{theorem}[No single horizon determines the Monoturn]
  \label{thm:horizon-nonunique}
  \statusproven{}
  Let $H_i \subsetneq \mathcal{M}$.
  There exist distinct global structures $\mathcal{M}_1 \neq \mathcal{M}_2$
  such that $F_i(\mathcal{M}_1) = F_i(\mathcal{M}_2)$.
  No reconstruction depending only on $F_i$ can distinguish them.
\end{theorem}

\begin{proof}
  Since $H_i \subsetneq \mathcal{M}$, there is a region
  $K = \mathcal{M} \setminus H_i \neq \emptyset$.
  Construct $\mathcal{M}_1$ and $\mathcal{M}_2$ agreeing on
  $H_i$ but differing on $K$:
  $\mathcal{M}_1|_{H_i} = \mathcal{M}_2|_{H_i}$,
  $\mathcal{M}_1|_K \neq \mathcal{M}_2|_K$.
  Then $F_i(\mathcal{M}_1) = F_i(\mathcal{M}_2)$ since $F_i$
  restricts to $H_i$.
  Any reconstruction $G_i \circ F_i$ assigns the same value
  to both, so it cannot identify which global structure
  is the actual one.
\end{proof}

% ---- Figures ------------------------------------------------

\[
  \mathcal{M}
  \xrightarrow{F_1} \text{Clusters}
  \xrightarrow{F_2} \text{Galaxies}
  \xrightarrow{F_3} \text{Atoms}
  \xrightarrow{F_4} \text{Fields}
  \xrightarrow{F_5} \text{Observer}
\]
\[
  \text{(Forward MSO: decomposition loses information at every arrow)}
\]

\medskip

\[
  \text{Observer}
  \xrightarrow{G_5} \text{Fields}
  \xrightarrow{G_4} \text{Atoms}
  \xrightarrow{G_3} \text{Galaxies}
  \xrightarrow{G_2} \text{Clusters}
  \xrightarrow{G_1} \widehat{\mathcal{M}}
\]
\[
  \text{(Reverse MSO: reconstruction recovers only what survived)}
\]

\begin{CitationAnchor}
  Mereology as a formal study of part-whole relations was
  initiated by Le\'sniewski (1916) and developed by Simons
  (1987) and Varzi.
  The present framework differs from classical mereology by
  treating decomposition and reconstruction as adjoint
  categorical operations ($F \dashv G$) rather than set-theoretic
  relations.
  Whitehead's process philosophy provides informal motivation
  for treating wholes as primary; see \textit{Process and Reality}
  (1929).
  The categorical formalisation follows Mac Lane~\cite{maclane1998}.
\end{CitationAnchor}

\begin{FurtherReading}
  Classical mereology is surveyed in Simons'
  \textit{Parts: A Study in Ontology} (1987).
  For a modern categorical perspective on part-whole relations,
  see Baez \& Stay (2011) and Spivak (2014).
\end{FurtherReading}
